\documentclass[14pt]{extarticle}
\usepackage{knu-fcc-style}

\begin{document}

\KNUTitlePageMake[
  type={Курсова робота},
  specialty={122 Комп'ютерні науки},
  title={Демонстрація елементів оформлення роботи за методичкою},
  department={Кафедра теоретичної кібернетики},
  studentInfo={Виконав студент 1-го курсу магістратури\\ОПП «Комп'ютерні науки»},
  student={Приклад Петро Іванович},
  studentSignature={false},
  supervisorInfo={Науковий керівник:\\доцент, к. ф.-м. н.},
  supervisor={Прикладенко Олена Сергіївна},
  supervisorSignature={false},
  declaration={Засвідчую, що в цій роботі\\немає запозичень з праць інших авторів\\без відповідних посилань.},
  declarationSigner={Студент},
  year={2026}
]

\KNUUnnumberedSection{Реферат}
\KNUWorkStats

\KNUKeywords{АЛГОРИТМ, ВЕРИФІКАЦІЯ, ДОКУМЕНТ, ОФОРМЛЕННЯ, СТИЛЬ}

Цей файл демонструє основні структурні елементи та приклади оформлення, описані в методичних вказівках факультету. Об'єктом демонстрації є шаблон навчальної роботи. Мета --- показати в одному документі титульний аркуш, реферат, зміст, перелік скорочень, вступ, розділи, підрозділи, рисунки, таблиці, переліки, примітки, виноски, формули, посилання, висновки, перелік джерел та додатки.

\KNUContents

\KNUUnnumberedSection{Перелік скорочень та умовних познак}
\noindent API --- програмний інтерфейс застосунку.\\
\noindent PDF --- Portable Document Format.\\
\noindent UI --- інтерфейс користувача.

\KNUUnnumberedSection{Вступ}
У вступі коротко пояснюють актуальність теми, мету, завдання, об'єкт і предмет дослідження. Цей демонстраційний документ показує, як можуть виглядати основні частини роботи, оформленої згідно з методичкою. У тексті також наведено приклад виноски\footnote{Приклад оформлення виноски безпосередньо в тексті.} та приклад посилання на джерело~\cite{knu-guide}.

\section{Структурні елементи та заголовки}
Цей розділ демонструє нумерацію розділів і правила оформлення заголовків.

\subsection{Підрозділ}
Підрозділи друкуються напівжирним шрифтом із абзацного відступу. У тексті можна робити посилання на рисунок~\ref{fig:demo}, таблицю~\ref{tab:demo}, формулу~\eqref{eq:demo} та додаток~\ref{app:extra}.

\subsubsection{Підпункт нижчого рівня}
Підпункти й подальші рівні також можуть використовуватися, якщо цього потребує структура роботи.

\paragraph{Приклад абзацного заголовка}
Текст після заголовка абзацного рівня має показувати, як виглядає ще один рівень структуризації матеріалу.

\section{Приклади оформлення елементів}
\subsection{Переліки}
Нижче наведено приклад простого переліку:
\begin{itemize}
  \item перший пункт переліку;
  \item другий пункт переліку;
  \item третій пункт переліку.
\end{itemize}

Нижче наведено приклад нумерованого переліку з вкладенням:
\begin{enumerate}
  \item перший пункт верхнього рівня;
  \item другий пункт верхнього рівня:
  \begin{enumerate}
    \item перший пункт другого рівня;
    \item другий пункт другого рівня.
  \end{enumerate}
\end{enumerate}

\subsection{Рисунок}
\begin{figure}[htbp]
  \centering
  \fbox{\rule{0pt}{45mm}\rule{0.72\textwidth}{0pt}}
  \caption{Приклад оформлення рисунка}
  \label{fig:demo}
\end{figure}

\subsection{Таблиця}
\begin{table}
  \caption{Приклад оформлення таблиці}
  \label{tab:demo}
  \centering
  \begin{tabular}{|l|c|c|}
    \hline
    Показник & Варіант 1 & Варіант 2 \\
    \hline
    Точність, \% & 91{,}5 & 94{,}2 \\
    \hline
    Час, с & 12{,}4 & 10{,}8 \\
    \hline
  \end{tabular}
\end{table}

\subsection{Примітка}
\vspace{2\baselineskip}
{\singlespacing\fontsize{12pt}{14pt}\selectfont\noindent\textbf{Примітка.} Це приклад одиночної примітки, поданої після основного тексту згідно з методичкою.\par}
\vspace{2\baselineskip}

\subsection{Приклад}
\begin{KNUExample}
Нехай потрібно показати окремо оформлений приклад із підписом «Приклад». Такий блок можна використовувати для демонстрації фрагмента тексту, формули або алгоритму.
\end{KNUExample}

\subsection{Формула}
Відомо, що
\begin{equation}
  \label{eq:demo}
  S = \sum_{i=1}^{n} i = \frac{n(n+1)}{2}.
\end{equation}
\noindent де $S$ --- сума перших $n$ натуральних чисел; $n$ --- кількість доданків.

\subsection{Фізичні одиниці та числа}
У науковому тексті важливо послідовно записувати числа та одиниці виміру: 12{,}5~МБ/с, 21{,}7~$^{\circ}$C, 3{,}14, 1~024 елементи.

\subsection{Посилання}
Приклад внутрішнього посилання: див. рисунок~\ref{fig:demo}, таблицю~\ref{tab:demo}, формулу~\eqref{eq:demo} та додаток~\ref{app:extra}. Приклад бібліографічного посилання: оформлення джерел описане в~\cite{dstu-style}.

\KNUUnnumberedSection{Висновки}
У висновках стисло підсумовують результати роботи. Для цього демонстраційного документа висновок полягає в тому, що один файл може показати більшість структурних елементів, описаних у методичці.

\begin{thebibliography}{9}
\bibitem{knu-guide} Методичні вказівки з підготовки та оформлення кваліфікаційних та курсових робіт для студентів факультету комп'ютерних наук та кібернетики. Київ, 2025.
\bibitem{dstu-style} ДСТУ 3008:2015. Інформація та документація. Звіти у сфері науки і техніки. Структура та правила оформлювання.
\bibitem{latex-companion} Mittelbach F., Goossens M. The \LaTeX{} Companion. 2nd ed. Boston: Addison-Wesley, 2004.
\end{thebibliography}

\appendix
\section{Додатковий матеріал}
\label{app:extra}
Цей додаток демонструє оформлення додатків. Тут можуть розміщуватися великі таблиці, додаткові рисунки, фрагменти коду або допоміжні результати.

\section{Ще один додаток}
Другий додаток показує, що додатки можуть бути кількома та нумеруються літерами.

\end{document}
